\chapter{Perpendiculaires et obliques}

\begin{itemize}
\it On considère un quadrilatère convexe $ABCD$ dans lequel les angles en $C$ et $D$ sont droits : 
\begin{enumerate}
\item Comparer $AD$ à $AC$ puis $AC$ à $AB+BC$.
\item En déduire que $AD<AB+BC$.
\end{enumerate}
\it Soit un triangle $ABC$ rectangle en $A$. La bissectrice intérieure de l'angle en $B$ coupe $(AC)$ en $D$, et on mène $(DE)$ perpendiculaire en $E$ à $(BC)$.
\begin{enumerate}
\item Comparer $AD$ et $DE$ puis $DE$ et $DC$. 
\item En déduire l'inégalité $AD<DC$. 
\end{enumerate}
\it On considère un triangle $ABC$ dans lequel on mène la hauteur $[AH]$.\begin{enumerate}
\item Comparer $AH$ à $AB$ puis à $AC$ et montrer que $2AH<AB+AC$.
\item En déduire que la somme des hauteurs est inférieure au périmètre du triangle $ABC$.
\end{enumerate}
\it Soit un triangle $ABC$ dans lequel on a $AB<AC$. Soit $[AM]$ la médiane issue de $A$. 
\begin{enumerate}
\item Démontrer que $\widehat{AMB}<\widehat{AMC}$ et que $\widehat{AMB}<90^o<\widehat{AMC}$.
\item En déduire que la hauteur $(AH)$ est du même côté que $(AB)$ par rapport à $(AM)$.
\end{enumerate}
\it On mène la bissectrice $[Oz)$ de l'angle $\widehat{xOy}$ et on porte deux longueurs égales 
$OA$ sur $[Ox)$ et $OB$ sur $[Oy)$.
\begin{enumerate}
\item Que représente $(Oz)$ pour le segment $[AB]$ ?
\item Montrer que pour tout point $M$ intérieur à l'angle $\widehat{xOz}$ on a $MA<MB$.
\end{enumerate}
\it Deux triangles isocèles $OAB$ et $OA'B'$ sont tels que $OA=OB=OA'=OB'$. Soient $H$ et $H'$ les milieux des bases $[AB]$ et $[A'B']$.
\begin{enumerate}
\item On suppose $\widehat{AOB}<\widehat{A'OB'}$. comparer $AB$ et $A'B'$ puis $AH$ et $A'H'$.
\item On suppose $AH<A'H'$. Comparer les angles $\widehat{AOB}$ et $\widehat{A'O'B'}$.
\item En déduire que : \emph{Si deux triangles rectangles ont même hypoténuse, les côtés de l'angle droit sont dans le même ordre de grandeur que les angles opposés.}
\end{enumerate}
\it Deux triangles rectangles $ABC$ et $A'BC$ sont situés du même côté de leur hypoténuse commune $[BC]$. 
On suppose $\widehat{BCA}<\widehat{BCA'}$ si bien que la droite $(CA)$ coupe le côté $[BA']$ en un point $D$ situé entre $B$ et $A'$.\begin{enumerate}
\item Montrer que l'angle $\widehat{CDB}$ est obtus, extérieur au triangle rectangle $\widehat{BAD}$ et que $D$ se trouve également entre $A$ et $C$.
\item En déduire que $BA<BA'$, $CA>CA'$ et $\widehat{ABC}>\widehat{A'BC}$. Étudier les réciproques.
\end{enumerate}
\it D'un point $A$, on mène la perpendiculaire $(AH)$ et les obliques $(AB)$ et $(AC)$ à la droite $(BC)$. \begin{enumerate}
\item Démontrer que l'une des égalités $HB=HC$, $AB=AC$, $\widehat{HAB}=\widehat{HAC}$, et $\widehat{HBA}=\widehat{HCA}$ entraîne les trois autres.
\item En supposant $B$ et $C$ d'un même côté de $H$ et en utilisant le théorème établi à l'exercice \textbf{50}, démontrer que l'une des inégalités $HB<HC$, $AB<AC$ , $\widehat{HAB}<\widehat{HAC}$ et $\widehat{HBA}>\widehat{HCA}$ entraîne les trois autres.
\end{enumerate}
\it On considère un angle $\widehat{xOy}$ et sa bissectrice $[Oz)$. Soit $M$ un point quelconque intérieur à l'angle $\widehat{xOz}$, $N$ et $P$ ses symétriques par rapport à $(Ox)$ et $(Oy)$, $A$ et $B$ ses projections sur $(Ox)$ et $(Oy)$. 
\begin{enumerate}
\item Comparer les triangles isocèles $MON$ et $MOP$ et en déduire que $MA<MB$.
\item Étudier réciproquement la région où se trouve le point $M$ intérieur à l'angle $\widehat{xOy}$ si on a par hypothèse $MA<MB$.
\end{enumerate}
\it On considère deux droites $(AA')$ et $(BB')$ se coupant en $O$ et les bissectrices $[Ox)$ et $[Oy)$ des angles $\widehat{AOB}$ et $\widehat{AOB'}$, ainsi que $[Ox')$ et $[Oy')$ des angles $\widehat{A'OB'}$ et $\widehat{A'OB}$.
\\ Démontrer en utilisant les résultats des exercices \textbf{58} ou \textbf{61} que : 
\begin{enumerate}
\item Tout point $M$ situé à l'intérieur de l'un des angles droits $\widehat{xOy}$ ou $\widehat{x'Oy'}$ est plus près de la droite $(AA')$ que de la droite $(BB')$.
\item Étudier et énoncer la réciproque de cette propriété.
\end{enumerate}
\end{itemize}